\documentclass[11pt,a4paper]{article}
\usepackage[T1]{fontenc}
\usepackage{lmodern}
\usepackage{microtype}
\usepackage[margin=25mm]{geometry}
\usepackage{amsmath,amssymb,amsthm,mathtools}
\usepackage{booktabs,array,longtable}
\usepackage{xcolor}
\usepackage{fancyvrb}
\usepackage[colorlinks=true,linkcolor=blue!40!black,urlcolor=blue!50!black,citecolor=blue!40!black]{hyperref}
\hypersetup{pdftitle={A rational five-cycle for two-block ADMM with multiplier step below two},pdfauthor={Dev D. Goyal}}
\newtheorem{theorem}{Theorem}
\newtheorem{lemma}[theorem]{Lemma}
\newtheorem{proposition}[theorem]{Proposition}
\newtheorem{corollary}[theorem]{Corollary}
\theoremstyle{remark}
\newtheorem{remark}[theorem]{Remark}
\newcommand{\R}{\mathbb R}
\newcommand{\conv}{\operatorname{conv}}
\newcommand{\ip}[2]{\langle #1,#2\rangle}
\newcommand{\norm}[1]{\lVert #1\rVert}
\newcommand{\argmin}{\operatorname*{arg\,min}}
\newcommand{\phig}{\frac{1+\sqrt5}{2}}
\setlength{\emergencystretch}{2em}
\setlength{\parskip}{3pt}
\title{\textbf{A rational five-cycle for two-block ADMM\\with multiplier step below two}}
\author{Dev D. Goyal}
\date{28 October 2026}
\begin{document}
\maketitle
\noindent\textbf{Autonomous AI research disclosure.} This manuscript and its proof sources were produced through fully autonomous agentic research. AI agents checked the supplied mathematical arguments, reran the exact arithmetic checks, and prepared this manuscript for release. The original model history is not documented in this release. No human proof review or external peer review is claimed. Dev D. Goyal is responsible for publication.\par\medskip
\begin{abstract}
Two compact convex polytopes in $\R^4$ are specified by explicit integer vertices. For these sets, classical two-block ADMM, with penalty one and multiplier-only step $199/100$, has a nonconstant period-five trajectory. Both primal subproblems have unique solutions, the polytopes contain the origin, and the instance has a primal-dual saddle point. The squared primal residual is at least $500000$ at every iteration. The construction refutes universal convergence for this exact update rule throughout $(0,2)$. Its relation to the question commonly associated with Glowinski is discussed with explicit literature limits. The proof consists of forty strict rational normal inequalities and an exact recurrence. Complete algebraic extensions are also given to finite strongly convex objectives and to $C^2$ strongly convex objectives with globally Lipschitz gradients and a unique constrained minimizer. Formal verification coverage and the limits of the literature comparison are stated separately.
\end{abstract}
\noindent\textbf{Verification status.} All seven local proof modules were recompiled with Lean 4.19.0 for this release. Eleven theorem endpoints passed the axiom audit. Their only axiom dependencies are \texttt{propext}, \texttt{Classical.choice}, and \texttt{Quot.sound}. The three exact arithmetic checkers also passed. Formal coverage is stated in Section~\ref{sec:formal}. The literature comparison remains provisional. No external peer review is claimed.


\section{The question and the exact algorithm}
For proper closed convex functions $f,g:\R^d\to\R\cup\{+\infty\}$, consider
\begin{equation}\label{eq:problem}
\min_{x,z\in\R^d} f(x)+g(z)\quad\text{subject to }x-z=0.
\end{equation}
With penalty parameter one, scaled multiplier $u$, and multiplier step $\tau$, the algorithm is
\begin{align}
x^{k+1}&=\argmin_x\left\{f(x)+\tfrac12\norm{x-z^k+u^k}^2\right\},\label{eq:admmx}\\
z^{k+1}&=\argmin_z\left\{g(z)+\tfrac12\norm{x^{k+1}-z+u^k}^2\right\},\label{eq:admmz}\\
u^{k+1}&=u^k+\tau(x^{k+1}-z^{k+1}).\label{eq:admmu}
\end{align}
All norms and inner products in this paper are Euclidean. Only the last update is multiplied by $\tau$.

The classical sufficient range for this multiplier-only step is
$0<\tau<\phig$. The question of extending the range to $(0,2)$ is
discussed by Tao and Yuan~\cite{TaoYuan}. Gu and Yang study the dual
step using performance estimation~\cite{GuYang}. These references supply
the research context. A complete primary-source literature comparison has not been established
for this release. No literature-priority claim is made.

The update rule is essential to the result. Generalized ADMM overrelaxation
also changes a primal argument. Relaxation of a Douglas--Rachford operator
can introduce an intermediate multiplier update. The present proof concerns
only \eqref{eq:admmx}--\eqref{eq:admmu}. Ryu and Yin give background on
these distinctions~\cite{RyuYin}. A convergence theorem for another update
rule does not imply convergence for the displayed rule.

\section{An explicit four-dimensional instance}
All subscripts on the following vectors are taken modulo five. Set $\tau=199/100$ and let rows of $Z,V$ be $z_i,v_i$ respectively, numbered from zero:
\begin{equation}\label{eq:dataZ}
Z=\begin{pmatrix}
20&-62&-161&80\\-52&-38&177&31\\-53&38&-126&-129\\20&62&26&178\\65&0&84&-160
\end{pmatrix},\qquad
V=\begin{pmatrix}
57&-289&314&-85\\-257&-143&-304&-116\\-216&200&178&273\\124&267&16&-325\\292&-35&-204&253
\end{pmatrix}.
\end{equation}
Define $u_i=\tau v_i$ and $x_i=z_i+v_i-v_{i-1}$. Thus the matrix with rows $x_i$ is
\begin{equation}\label{eq:dataX}
X=\begin{pmatrix}
-215&-316&357&-258\\-366&108&-441&0\\-12&381&356&260\\360&129&-136&-420\\233&-302&-136&418
\end{pmatrix}.
\end{equation}
Put
\begin{equation}\label{eq:sets}
C=\conv\{x_0,x_1,x_2,x_3,x_4\},\qquad
D=\conv\{z_0,z_1,z_2,z_3,z_4\}.
\end{equation}
For a set $S$, its convex indicator is $\delta_S(w)=0$ for $w\in S$ and $+\infty$ otherwise.

\begin{theorem}[Nonconvergence below two]\label{thm:main}
Apply \eqref{eq:admmx}--\eqref{eq:admmu} to $f=\delta_C$, $g=\delta_D$, with penalty one, $\tau=199/100$, and initial state $(z^0,u^0)=(z_0,u_0)$. Every primal subproblem has a unique minimizer. For every $k\ge0$, with the harmless initialization $x^0=x_0$,
\[
(x^k,z^k,u^k)=(x_{k\bmod5},z_{k\bmod5},u_{k\bmod5}).
\]
The sets are compact and convex, contain zero, and the problem has the saddle point $(0,0,0)$. Nevertheless the state is nonconstant and five-periodic, and
\[
\norm{x^k-z^k}^2\ge 500000\quad\text{for all }k\ge0.
\]
In particular, the primal residual does not tend to zero. Hence classical multiplier-only two-block ADMM does not converge universally throughout $(0,2)$.
\end{theorem}

\subsection{The finite certificate}
Define
\begin{equation}\label{eq:normals}
a_i=z_{i-1}-u_{i-1}-x_i,\qquad b_i=x_i+u_{i-1}-z_i.
\end{equation}
The rows of the following matrices are $100a_i$ and $100b_i$:
\begin{equation}\label{eq:normalA}
100A=\begin{pmatrix}
-30108&38565&13296&-40547\\27257&40511&-34486&24915\\47143&-13443&42596&184\\1684&-48900&-34422&-25227\\-45976&-16733&13016&40675
\end{pmatrix},
\end{equation}
\begin{equation}\label{eq:normalB}
100B=\begin{pmatrix}
34608&-32365&11204&16547\\-20057&-42911&686&-20015\\-47043&5843&-12296&15816\\-8984&46500&19222&-5473\\41476&22933&-18816&-6875
\end{pmatrix}.
\end{equation}
Here $A,B$ are only names for the displayed normal matrices. The constraint matrices in \eqref{eq:problem} remain $I,-I$.

Direct integer arithmetic gives all twenty-five entries of each normal table, including its five zero diagonal entries:
\begin{equation}\label{eq:marginC}
\left(100\ip{a_i}{x_i-x_j}\right)_{i,j=0}^4=
\begin{pmatrix}
0&173466&248761&136989&46913174\\
47008755&0&298961&343217&386327\\
252974&47015535&0&157520&196556\\
264620&292314&47039232&0&278194\\
247902&292936&187297&47136525&0
\end{pmatrix},
\end{equation}
\begin{equation}\label{eq:marginD}
\left(100\ip{b_i}{z_i-z_j}\right)_{i,j=0}^4=
\begin{pmatrix}
0&292387&8829067&296506&1675570\\
1626843&0&246637&8780995&218220\\
282535&1686239&0&307387&8845564\\
8824160&296099&1701701&0&322530\\
328346&8787159&284479&1676996&0
\end{pmatrix}.
\end{equation}
There are forty off-diagonal inequalities in total, and all are strict. The minimum margin is $136989/100$. In addition,
\begin{equation}\label{eq:centroids}
\sum_i x_i=\sum_i z_i=\sum_i v_i=0,
\end{equation}
and the five squared residuals are, in index order,
\begin{equation}\label{eq:residuals}
(\norm{x_i-z_i}^2)_{i=0}^4=(502309,502797,502975,503937,501912).
\end{equation}
Every quantity in \eqref{eq:normalA}--\eqref{eq:residuals} is independently reproducible from the two integer arrays \eqref{eq:dataZ}.

\subsection{Proof over the entire convex sets}
The transition from the finite tables to the actual minimizers is elementary but essential.
\begin{lemma}[Projection certificate]\label{lem:projection}
Let $S=\conv\{p_0,\ldots,p_m\}\subset\R^d$. If $p=p_i$ and $q\in\R^d$ satisfy
\[
\ip{q-p}{p_j-p}\le0\quad\text{for every }j,
\]
then $p$ is the unique minimizer of $\norm{q-y}^2$ over all $y\in S$.
\end{lemma}
\begin{proof}
For $y=\sum_j\lambda_j p_j\in S$, where $\lambda_j\ge0$ and $\sum_j\lambda_j=1$, linearity gives
\[
\ip{q-p}{y-p}=\sum_j\lambda_j\ip{q-p}{p_j-p}\le0.
\]
Consequently
\begin{equation}\label{eq:projectionidentity}
\norm{q-y}^2-\norm{q-p}^2
=\norm{y-p}^2-2\ip{q-p}{y-p}\ge\norm{y-p}^2.
\end{equation}
This proves minimality, and equality can hold only at $y=p$.
\end{proof}

\begin{proof}[Proof of Theorem~\ref{thm:main}]
Finite convex hulls are nonempty, compact, and convex, hence closed. Equation~\eqref{eq:centroids} expresses zero as the uniform average of the vertices of each hull, so $0\in C\cap D$. Thus the primal problem has optimal value zero. For its unaugmented Lagrangian
\[
L(x,z,u)=\delta_C(x)+\delta_D(z)+\ip{u}{x-z},
\]
one has $L(0,0,u)=0=L(0,0,0)\le L(x,z,0)$ for all $u,x,z$, in the extended-real sense. This proves the stated saddle-point property.

Fix $i$. By \eqref{eq:marginC}, $\ip{a_i}{x_j-x_i}\le0$ for every $j$. Applying Lemma~\ref{lem:projection} with $q=z_{i-1}-u_{i-1}$ proves
\begin{equation}\label{eq:projX}
x_i=P_C(z_{i-1}-u_{i-1}).
\end{equation}
The same argument using \eqref{eq:marginD} proves
\begin{equation}\label{eq:projZ}
z_i=P_D(x_i+u_{i-1}).
\end{equation}
These are the unique minimizers in \eqref{eq:admmx} and \eqref{eq:admmz} for the indicator objectives, since points outside the respective sets have infinite objective value. Compactness and strict convexity of the squared distance also guarantee existence and uniqueness for every possible projection center.

By the definition of $x_i$,
\begin{equation}\label{eq:dualidentity}
u_{i-1}+\tau(x_i-z_i)
=\tau v_{i-1}+\tau(v_i-v_{i-1})
=\tau v_i=u_i.
\end{equation}
Equations~\eqref{eq:projX}--\eqref{eq:dualidentity} therefore give one complete ADMM update from phase $i-1$ to phase $i$. Induction, starting at phase zero, proves the asserted orbit for every natural-number iteration. Its five $z_i$ are distinct, so it is nonconstant and has period five. Finally \eqref{eq:residuals} gives the uniform lower bound. A real sequence bounded below by $500000$ cannot tend to zero.
\end{proof}

\begin{remark}[Nondegenerate feasible geometry]
The four-row affine determinants, using row four as base point, are
\[
\det\begin{pmatrix}x_0-x_4\\x_1-x_4\\x_2-x_4\\x_3-x_4\end{pmatrix}=395537260080,
\qquad
\det\begin{pmatrix}z_0-z_4\\z_1-z_4\\z_2-z_4\\z_3-z_4\end{pmatrix}=1913707790.
\]
Both sets are thus full-dimensional simplices. Together with the strictly positive uniform barycentric coefficients in \eqref{eq:centroids}, this puts zero in the ordinary interior of each. These determinant and interior statements are exact auxiliary checks; the Lean-verified core does not require them.
\end{remark}

\section{Extensions with complete algebraic proofs}
The core result is Theorem~\ref{thm:main}. The following consequences give additional information about the mechanism. Formal coverage is recorded in Section~\ref{sec:formal}; properties not included in a checked Lean declaration are not represented as formally verified.

\subsection{A whole interval for the same polytopes}
Keep $X,Z,V$ fixed and define $u_i=\tau v_i$ for a variable $\tau$. Every positive margin is affine in $\tau$:
\begin{align*}
m^C_{ij}(\tau)&=\ip{z_{i-1}-x_i}{x_i-x_j}
-\tau\ip{v_{i-1}}{x_i-x_j},\\
m^D_{ij}(\tau)&=\ip{x_i-z_i}{z_i-z_j}
+\tau\ip{v_{i-1}}{z_i-z_j}.
\end{align*}
For $i\ne j$ all forty slopes are strictly positive. Computing their forty roots in exact rational arithmetic gives
\begin{equation}\label{eq:threshold}
\max_{i\ne j}\{\text{root}(m^C_{ij}),\text{root}(m^D_{ij})\}
=\frac{140111}{70637}.
\end{equation}
The binding entry is $(i,j)=(0,3)$ in the $C$ table, for which
\[
m^C_{03}(\tau)=-420333+211911\tau.
\]
The supplied exact verifier includes every slope and root, and checks \eqref{eq:threshold} by integer comparisons. Thus the same two polytopes have the displayed cycle for every $\tau\ge140111/70637$. Non-strict normal inequalities at the endpoint still give unique projections by \eqref{eq:projectionidentity}. In particular, the counterexample covers the interval $[140111/70637,2)$, whose lower endpoint is approximately $1.98353554$. No claim of sharpness is made for this threshold or for the golden ratio.

\subsection{Finite strongly convex objectives}
For the remainder of this section, return to $\tau=199/100$. Let $\mu=1/1000$ and set
\[
q_i=a_i-\mu x_i,\qquad r_i=b_i-\mu z_i.
\]
The same integer data give the strict bounds
\begin{align}
\ip{q_i}{x_i-x_j},\ \ip{r_i}{z_i-z_j}
&\ge \frac{24286}{25}>1 &&(i\ne j),\label{eq:strongcross}\\
\ip{q_i}{x_i},\ \ip{r_i}{z_i}
&\ge \frac{2713619}{125}>1 &&(0\le i\le4).\label{eq:strongorigin}
\end{align}
These rational inequalities are checked directly in the standalone verifier; they can also be obtained by subtracting the indicated quadratic terms from the normal tables.

Define
\begin{align}
f_{\rm plq}(w)&=\frac\mu2\norm{w}^2+
\max\left(0,\max_i\ip{q_i}{w-x_i}\right),\label{eq:plqf}\\
g_{\rm plq}(w)&=\frac\mu2\norm{w}^2+
\max\left(0,\max_i\ip{r_i}{w-z_i}\right).\label{eq:plqg}
\end{align}
They are finite, continuous, and $\mu$-strongly convex. Their nonquadratic terms vanish both at zero and at their respective sampled points. The affine term indexed by $i$ is active at $x_i$. Therefore for every $w$,
\begin{align}
f_{\rm plq}(w)-f_{\rm plq}(x_i)
&\ge \frac\mu2(\norm{w}^2-\norm{x_i}^2)+\ip{q_i}{w-x_i}\notag\\
&=\ip{a_i}{w-x_i}+\frac\mu2\norm{w-x_i}^2.\label{eq:strongsupport}
\end{align}
The analogous inequality holds for $g_{\rm plq}$ with normal $b_i$. Combining \eqref{eq:strongsupport} with the quadratic identity \eqref{eq:projectionidentity} proves that each ADMM subproblem has exactly the same unique minimizer as before: the objective difference from that minimizer is at least $(1+\mu)\norm{w-x_i}^2/2$, or its $z$ analogue. Equation~\eqref{eq:dualidentity} then proves the same infinite orbit.

Both functions are bounded below by $\mu\norm{w}^2/2$ and vanish at zero. Thus under the constraint $x=z$, zero is the unique minimizer. At multiplier zero this is also a saddle point. The nonconvergence therefore persists with everywhere-finite strongly convex objectives and a unique solution.

\subsection{Explicit \texorpdfstring{$C^2$}{C2} strongly convex objectives}
Define a scalar function by
\begin{equation}\label{eq:psi}
\psi(t)=\begin{cases}
0,&t\le-1,\\
(t+1)^3-\tfrac12(t+1)^4,&-1\le t\le0,\\
t+\tfrac12,&t\ge0.
\end{cases}
\end{equation}
The formulas agree at both endpoints. On the middle interval put $s=t+1\in[0,1]$. Then
\[
\psi'(t)=3s^2-2s^3,\qquad
\psi''(t)=6s(1-s).
\]
The first two derivatives agree with the exterior formulas at both endpoints. Hence $\psi$ is $C^2$, convex, and satisfies $0\le\psi''\le3/2$ and $\psi'(0)=1$.

Set
\begin{equation}\label{eq:smoothf}
f_{\rm sm}(w)=\frac\mu2\norm{w}^2+
\sum_{i=0}^4\psi\bigl(\ip{q_i}{w-x_i}\bigr),
\end{equation}
\begin{equation}\label{eq:smoothg}
g_{\rm sm}(w)=\frac\mu2\norm{w}^2+
\sum_{i=0}^4\psi\bigl(\ip{r_i}{w-z_i}\bigr).
\end{equation}
These are $C^2$ and $\mu$-strongly convex. Their Hessians have the form
\[
\nabla^2 f_{\rm sm}(w)=\mu I+
\sum_i\psi''\bigl(\ip{q_i}{w-x_i}\bigr)q_iq_i^\mathsf{T}.
\]
Consequently a valid global Lipschitz constant for $\nabla f_{\rm sm}$ is $\mu+(3/2)\sum_i\norm{q_i}^2$, with the analogous bound for $g_{\rm sm}$.

At $x_j$, all hinge arguments except the $j$th are below $-1$ by \eqref{eq:strongcross}; the $j$th argument is zero. It follows that
\[
\nabla f_{\rm sm}(x_j)=\mu x_j+q_j=a_j,
\qquad \nabla g_{\rm sm}(z_j)=\mu z_j+r_j=b_j.
\]
These are exactly the first-order conditions for \eqref{eq:admmx} and \eqref{eq:admmz}. Strong convexity makes each minimizer unique. All hinge arguments at the origin are below $-1$ by \eqref{eq:strongorigin}, so both objectives vanish there. They are nonnegative and bounded below by the positive quadratic term. Thus zero is again the unique constrained optimizer, while the identical five-cycle persists.

There is also a purely algebraic certificate of the proximal updates. For every real $t$,
\begin{equation}\label{eq:hingesupport}
\psi(t)\ge0,\qquad \psi(t)\ge t+\tfrac12,\qquad \psi(0)=\tfrac12.
\end{equation}
On the middle interval these follow respectively from
\[
\psi(t)=\tfrac12s^3(2-s),\qquad
\psi(t)-(t+\tfrac12)=\tfrac12(1-s)^3(1+s).
\]
The exterior intervals are immediate. At a sampled point all other hinge terms vanish, so \eqref{eq:hingesupport} gives \eqref{eq:strongsupport} for the smooth objectives as well. This proves the same unique minimizers without relying on differentiation.

\section{Constructing a finite certificate}
The search used the standard Gram-matrix idea from performance estimation, which has already been applied to ADMM~\cite{GuYang}. For a proposed period $K$, take formal vectors $z_i,u_i$ and express
\[
x_i=z_i+\frac{u_i-u_{i-1}}{\tau},\quad
a_i=z_{i-1}-u_{i-1}-x_i,\quad b_i=x_i+u_{i-1}-z_i.
\]
All these vectors depend linearly on the $2K$ basic vectors. Their Gram matrix $G$ can therefore be chosen by a semidefinite program imposing
\[
G\succeq0,\quad \operatorname{tr}G=1,\quad
\sum_i z_i=\sum_i u_i=0,\quad
\ip{a_i}{x_j-x_i},\ \ip{b_i}{z_j-z_i}\le-\eta\quad(i\ne j),
\]
and maximizing $\eta$. Centering is imposed by requiring the corresponding coefficient vectors to lie in the kernel of $G$. A simultaneous cyclic permutation of the two families may also be imposed as a linear symmetry of $G$.

The indicator construction assigns the same function value to every sampled vertex. As a result, the pairwise supporting inequalities directly certify a single convex set for each block. A numerical rank-four solution with $K=5$ and $\tau=1.99$ supplied positive slack. After factoring its Gram matrix, the construction set $v_i=u_i/\tau$, scaled and rounded $z_i,v_i$, and chose the last row of each family to make its centroid exactly zero. The resulting integer arrays are \eqref{eq:dataZ}. Their acceptance depends entirely on the rational tables and the proof above. The semidefinite solver is used only for discovery; none of its output or feasibility tolerances is trusted in the final proof. Portable discovery scripts and the original numerical candidate are included for provenance.

\section{Formal verification and reproducibility}\label{sec:formal}
The formal core uses Lean 4.19.0 and mathlib v4.19.0. It works over actual real coordinate vectors, with explicitly defined Euclidean dot product and squared distance,
\[
\operatorname{dot}(p,q)=\sum_{j=0}^3p_jq_j,
\qquad\operatorname{sqDist}(p,q)=\sum_{j=0}^3(p_j-q_j)^2.
\]
The ambient function type's default norm is never substituted for these expressions. The feasible sets are the actual real convex hulls of the displayed five vertices.

\begin{center}
\begin{tabular}{@{}p{0.31\linewidth}p{0.63\linewidth}@{}}
\toprule
Module & Mathematical content\\\midrule
\texttt{ADMMGeometry.lean} & Convexity, compactness, closedness, zero-centroid membership; extension of the vertex inequalities to every point of a convex hull; unique projection theorem.\\[4pt]
\texttt{ADMMData.lean} & Explicit real vectors, cyclic arithmetic, multiplier identity, all normal inequalities, and the residual lower bound.\\[4pt]
\texttt{ADMMCounterexample.lean} & Every natural-number ADMM update, five-periodicity, nonconstancy, feasible zero, a zero fixed trajectory, and failure of residual convergence.\\[4pt]
\texttt{ADMMHingeScalar.lean},\newline\texttt{ADMMHinge.lean},\newline\texttt{ADMMHingeTrajectory.lean} & Algebraic scalar hinge bounds, the explicit finite objectives' global supporting inequalities, unique proximal minimizers, their unique constrained optimizer, and the infinite nonconvergent trajectory.\\\bottomrule
\end{tabular}
\end{center}

The formal projection predicate includes membership, global minimality over the whole convex set, and uniqueness. Its hypotheses are discharged from the explicit numerical data; the final theorem assumes neither the normal inequalities nor nonconvergence. The predicate \texttt{IsADMMTrajectory} requires, for every natural number $n$, the unique projections with centers $Z_n-U_n$ and $X_{n+1}+U_n$, followed by the multiplier identity $U_{n+1}=U_n+\tau(X_{n+1}-Z_{n+1})$. The checked source files are included in full in the reproducible artifact.

\noindent\textbf{Checked core.} The declaration \texttt{ADMMCounterexample.counterexample} compiled with exit code zero. It contains $1<\tau<2$, convexity, compactness, and closedness of both sets, the common feasible origin, a zero fixed trajectory, every update of the displayed trajectory, five-periodicity, nonconstancy, the uniform squared-residual bound, and the assertion that this residual does not tend to zero. The actual projection minimizers are checked over the full real convex hulls, rather than only at their vertices.

The three core modules have no \texttt{sorry}, \texttt{admit}, custom \texttt{axiom}, \texttt{native\_decide}, or externally trusted numerical certificate. The final axiom report is exactly \texttt{[propext, Classical.choice, Quot.sound]}. These are the standard logical foundations, not assumptions of the counterexample. The use of \texttt{noncomputable} for real arithmetic changes executable-code generation, not the proof obligation or the kernel's checking of the theorem.

The indicator-objective and saddle-point interpretation is proved in the text. The affine determinants, ordinary interiority, the parameter interval, and analytic $C^2$/strong-convexity/Lipschitz-gradient properties of the extensions are not part of this core Lean declaration. Their complete mathematical arguments and exact arithmetic certificates are supplied separately; no formal-verification claim is made for those additional regularity statements.

\noindent\textbf{Checked finite-objective extension.} The separate declaration \path{ADMMHingeTrajectory.verified_orbit_and_optimizer} also compiled successfully, with the same three standard foundational dependencies. For the actual functions in \eqref{eq:smoothf}--\eqref{eq:smoothg}, it certifies global proximal minimizers at every iteration, the multiplier recurrence, the unique constrained optimizer zero, periodicity, nonconstancy, and the failure of the residual to tend to zero. The analytic regularity statements about these functions are established by the handwritten derivative and Hessian calculations above; those properties are not encoded in this additional Lean theorem.


The exact Python certificate uses only the standard library and rational arithmetic. Two separate arithmetic checkers reconstruct the data, normals, determinants, residuals, and parameter threshold. Discovery used semidefinite convex interpolation followed by rational reconstruction. The numerical solver and its floating-point output are not trusted by the proof or required to reproduce it.

The bundle includes a pinned Lake project and a direct compiler runner. The latter records the compiler version, input hashes, exit codes, and axiom reports. Both routes are described in its README. The release verification receipt records the source hashes, module results, and axiom dependencies from the fresh replay.

\section{Research context, implications, and limits}
The mathematical conclusion is independent of a priority assessment.
The displayed exact cycle refutes universal convergence of
\eqref{eq:admmx}--\eqref{eq:admmu} throughout $(0,2)$, even for the
explicit smooth strongly convex objectives. It leaves the usual unit-step
two-block ADMM theory intact.

Relevant references include Tao--Yuan~\cite{TaoYuan},
Gu--Yang~\cite{GuYang}, Chen--Sun--Toh~\cite{ChenSunToh},
Ma et al.~\cite{MaEtAl}, and He--Ma--Yuan~\cite{HeMaYuan}.
The literature comparison is provisional. Detailed primary-source comparisons
of special model assumptions, conditional convergence, proof-measure
obstructions, and alternative multiplier updates have not been independently
confirmed for this release. These references are retained for context and
attribution. Their inclusion does not establish novelty. This manuscript
makes no claim of first discovery or a complete literature review.

The construction does not determine the largest universally convergent
multiplier interval. In particular, failure at $199/100$, or on
\eqref{eq:threshold}, does not show that the golden ratio is the sharp
boundary. The orbit and global minimizer claims are formally checked;
the analytic regularity claims are proved in the text. Neither check
establishes literature priority or external expert validation.

{\small
\begin{thebibliography}{9}
\setlength{\itemsep}{1pt}
\setlength{\parskip}{0pt}
\bibitem{TaoYuan} M. Tao and X. Yuan. \emph{On Glowinski's Open Question on the Alternating Direction Method of Multipliers}. Journal of Optimization Theory and Applications 179 (2018), 163--196. \href{https://doi.org/10.1007/s10957-018-1338-x}{doi:10.1007/s10957-018-1338-x}. Primary preprint: \url{https://optimization-online.org/wp-content/uploads/2017/06/6084.pdf}.
\bibitem{GuYang} G. Gu and J. Yang. \emph{On the dual step length of the alternating direction method of multipliers}. arXiv:2006.08309v2 (2020). \url{https://arxiv.org/abs/2006.08309}.
\bibitem{ChenSunToh} L. Chen, D. Sun, and K.-C. Toh. \emph{A Note on the Convergence of ADMM for Linearly Constrained Convex Optimization Problems}. Computational Optimization and Applications 66 (2017), 327--343. \href{https://doi.org/10.1007/s10589-016-9864-7}{doi:10.1007/s10589-016-9864-7}. See published Theorem 4.1 and Remark 4.3.
\bibitem{MaEtAl} F. Ma, J. Shu, Y. Li, and J. Wu. \emph{The dual step size of the alternating direction method can be larger than 1.618 when one function is strongly convex}. Journal of Industrial and Management Optimization 17 (2021), 1173--1185. \href{https://doi.org/10.3934/jimo.2020016}{doi:10.3934/jimo.2020016}.
\bibitem{RyuYin} E. K. Ryu and W. Yin. \emph{Large-Scale Convex Optimization: Algorithms and Analyses via Monotone Operators}. Cambridge University Press, 2022. Author PDF compiled 20 April 2023, Chapter 8 and Exercise 8.1. \url{https://large-scale-book.mathopt.com/LSCOMO.pdf}.
\bibitem{Zhang} T. Zhang. \emph{On the $O(1/K^2)$ Ergodic Convergence of ADMM with Dual Step Size from 0 to 2}. Journal of the Operations Research Society of China 14 (2026), 521--532. \href{https://doi.org/10.1007/s40305-024-00568-7}{doi:10.1007/s40305-024-00568-7}.
\bibitem{XuWang} K. Xu and X. Wang. \emph{AI-Assisted Discovery and Construction of a Counterexample to the Convergence of Three-Block ADMM with the Identity Matrix as its Third Constraint Block}. arXiv:2608.14396 (2026). \url{https://arxiv.org/abs/2608.14396}.
\bibitem{AnshikaEtAl} A. Anshika, J. Li, D. Ghosh, and X. Zhang. \emph{A three-operator splitting scheme derived from three-block ADMM}. Optimization and Engineering 27 (2026), 2119--2159. \href{https://doi.org/10.1007/s11081-025-10060-7}{doi:10.1007/s11081-025-10060-7}. Published online 19 December 2025; Section 5.1.
\bibitem{HeMaYuan} B. He, F. Ma, and X. Yuan. \emph{On the Step Size of Symmetric Alternating Directions Method of Multipliers}. Primary preprint, 26 May 2015. \url{https://optimization-online.org/wp-content/uploads/2015/05/4925.pdf}. Published as \emph{Convergence Study on the Symmetric Version of ADMM with Larger Step Sizes}, SIAM Journal on Imaging Sciences 9 (2016), 1467--1501, \href{https://doi.org/10.1137/15M1044448}{doi:10.1137/15M1044448}.
\end{thebibliography}
}
\end{document}
